\section{同一个问题怎样跨核区比较}\label{sec:4-1}

在第~1~章中，原子核低能激发谱已被置于六个典型质量区的宏观坐标之中；第~2~章与第~3~章则进一步从微观壳模型、平均场自洽形变、对称性恢复投影以及转动、手征与摆动等集体动力学角度，奠定了多体波函数与唯象图像的数理基础。当分析视角从单个核素转向全局核素图谱时，跨核区比较往往成为检验模型泛化能力、揭示演化规律的最有力工具。然而，跨核区比较也是极易滋生虚假类比与概念滑移的区域。

最常见的认识论误区在于：仅凭某一核区发现了特定能带结构（如两条近简并带或交错低自旋带），便将质量数 $A$ 或单一唯象参数（如能级比 $E(2^+_2)/E(2^+_1)$、三轴形变角 $\gamma$ 的单一平衡值）作为排他性判据，孤立地宣布在另一区域也必定存在完全相同的物理模式。事实上，核多体系统的集体响应从来不是孤立几何形状的直接映射，而是价轨道占据、有效相互作用单极漂移、配对超流关联与自旋窗口内探针选择定则综合作用的宏观显现。本节旨在建立一套跨核区比较的物理判据与分析方法，明确哪些量子几何与多极代数是普适可迁移的，哪些微观壳结构与有效参数是严格局域不可迁移的。

\subsection{比较对象：内禀组态、实验室对称性与态的映射}\label{sec:4-1-1}

跨核区比较的首要前提是厘清``究竟在比较什么对象''。原子核作为自约束的有限量子多体系统，在没有外场作用的实验室参考系中，由强相互作用与电磁相互作用主导的哈密顿量在极高精度下严格具备空间旋转群 $\mathrm{SO}(3)$ 对称性、空间反演宇称 $\mathcal{P}$ 对称性以及时间反演 $\mathcal{T}$ 对称性（在此忽略微弱的弱相互作用宇称不守恒效应）。因此，实验室参考系中物理定态的基矢必定可以按确定的总角动量量子数 $I$、宇称 $\pi$ 及其磁量子数 $M$ 展开。在实际谱学实验中，测得的通常是具有确定能级位置的纯态或处于特定取向系综与统计混合的退极化布居。

自洽平均场模型中所讨论的轴对称长椭球、扁椭球或三轴椭球形状，本质上是内禀对称性破缺参考态 $|\Phi(\beta,\gamma)\rangle$ 的几何表象。内禀系中的非零四极形变参量 $(\beta, \gamma)$ 描述的是核内多核子空间密度分布的二阶各向异性关联，如第~2.4~节式~\eqref{eq:2-4-2}所引入的玻尔坐标。

然而，单一内禀态 $|\Phi(\beta,\gamma)\rangle$ 本身并非实验室物理本征态。利用角动量投影算符 $\hat{P}^I_{MK}$ 对单一内禀构型进行全欧拉角积分，所构造的仅是具备良好角动量量子数的投影试探态（trial states）：
\begin{equation}\label{eq:4-1-projection}
|IM\alpha\rangle = \sum_K g_K^\alpha \hat{P}^I_{MK} |\Phi\rangle = \sum_K g_K^\alpha \frac{2I+1}{8\pi^2} \int \mathrm{d}\Omega \, \mathcal{D}^{I*}_{MK}(\Omega) \hat{R}(\Omega) |\Phi\rangle.
\end{equation}
在物理真实的有限核系统中，实际定态通常是多个具有不同形变构型、配对自由度或准粒子激发的内禀态的多体量子叠加（例如生成坐标法 GCM 或多参考能量泛函理论所描述的构型混合网络）。

这意味着，跨核区比较的对象绝不能停留在某个未投影平均场势能面（PES）局域极小点的坐标数值上，而必须落实到经对称性恢复后具有确定自旋与宇称的物理态 $|I^\pi\rangle$、这些态之间的跃迁网络以及波函数在形变构型空间上的概率分布。在波函数分布极其弥散的大振幅软核中，实验室一体电荷密度对于 $0^+$ 基态处处呈现球对称，唯有通过高阶两体密度关联与电磁多极约化矩阵元，方能提取内部形状的客观物理实在性\cite{poves-2020-limits-assigning-nuclear-shape}。

\subsection{微观驱动与动力学响应：轨道占据、单极演化与配对关联}\label{sec:4-1-2}

不同质量区的最显著本质差异，在于费米面所处的主壳层尺度与高 $j$ 独特宇称轨道（intruder orbitals）的微观占据状态。例如：
\begin{itemize}
\item 在 $A\approx 60$ 与 $80$ 区域，质子与中子费米面主要受 $1f_{5/2}, 2p_{3/2}, 2p_{1/2}$ 以及下落的 $1g_{9/2}$ 独特宇称轨道主导；
\item 在 $A\approx 130$ 区域，中子接近壳层中上部，质子位于准空穴或下部，负宇称独特轨道 $1h_{11/2}$ 的驱动效应成为主导；
\item 在 $A\approx 160$ 稀土重区，强形变长椭球势阱将 $1i_{13/2}$ 轨道大幅劈裂；
\item 在 $A\approx 190$ 区域，费米面则深埋于 $1i_{13/2}$ 中子空穴与 $1h_{9/2}, 2f_{7/2}$ 混合轨道之中。
\end{itemize}

单粒子轨道对于核芯形变的偏好完全由其角动量在内禀对称轴上的投影 $K$（或 Nilsson 量子数 $\Omega$）决定。根据微观单粒子势能梯度，当费米面位于高 $j$ 独特轨道的底部（低 $\Omega$ 分量，如 $1/2^-, 3/2^-$）时，波函数沿赤道平面外延，粒子与核心耦合倾向于驱动核芯发生长椭形变（prolate）；当费米面位于高 $j$ 轨道的顶部（高 $\Omega$ 分量，如 $11/2^-, 9/2^-$）时，波函数偏向于对称轴极区，表现为空穴型行为，驱动核芯发生扁椭形变（oblate）；而当费米面恰处于高 $j$ 壳层中部时，粒子态与空穴态的横向耦合对长轴与短轴施加相互竞争的正交驱动力，从而成为诱发三轴不稳定性（triaxiality）的微观温床\cite{book-bohr-mottelson-nuclear-structure-vol2}。

更深层次的微观驱动还来自于多核子相互作用中的单极效应（monopole interaction）。正如第~2.2~节式~\eqref{eq:2-2-1}所展开的单极相互作用分解与有效单粒子能量（ESPE）演化框架所示，随着价核子在不同组态间的填充或跃迁，核子间张量力与自旋-同位旋交换相互作用引起有效能隙的动态漂移\cite{doctoral-r1-g28-otsuka-2023-type-ii-shell-evolution}。在跨区比较中，这种单极演化逻辑展现出鲜明的核区特性：例如在 $A\approx 100$ 区域，正是质子从 $2p, 1f$ 轨道激发至 $1g_{9/2}$ 时，通过强烈的单极吸引作用诱发了中子 $1g_{7/2}$ 与 $1h_{11/2}$ 轨道的突变性下移，形成了所谓的第二类壳演化（Type-II shell evolution）与形变自强化反馈\cite{doctoral-r1-g28-togashi-2016-zr-mcsm-qpt}；而在高质子数区，这种单极诱导相变的临界核子数则发生显著位移。

配对关联（pairing correlations）在跨区比较中构成了另一项不可忽缺的动力学调制。单极配对力平滑了单粒子费米面处的占据阶跃，使得准粒子玻戈留波夫变换成为必需。正如第~2.6.3~节在讨论冻结场的 Inglis--Belyaev 响应时所推导的，配对超流关联抑制了准粒子非对角激发，使得低自旋下的转动惯量显著低于刚体惯量值。然而必须指出，配对导致的惯量削减程度绝非普适常数：不同质量区由于单粒子能级密度 $g(\epsilon_{\mathrm{F}})$ 的巨大差异，配对能隙 $\Delta$ 呈现剧烈的系统性起伏；而当外加转动自旋升高引起科里奥利对拆裂（Coriolis anti-pairing effect）时，带交叉与对淬灭的速度在不同区域截然不同。因此，跨区比较能带动力学时，绝不能将某一区域拟合所得的有效转动惯量直接作为通用常数跨区移植。

\subsection{形变空间测度与高阶不变量：$\gamma$ 刚性、$\gamma$ 软性与分布宽度}\label{sec:4-1-3}

三轴自由度是跨区比较中最易引起混淆的概念之一。在唯象模型中，人们常将 Davydov-Filippov 刚性三轴转子模型与 Wilets-Jean $\gamma$ 不稳定模型进行对照：
\begin{itemize}
\item \textbf{刚性三轴假设}：设想势能面在非零三轴平衡角（如 $\gamma \approx 20^\circ\text{--}30^\circ$）处具有较深而局域的极小，波函数集中于极小点近旁；
\item \textbf{$\gamma$ 软性/不稳定假设}：势能面沿 $\gamma$ 方向展现宽平坦特性，核波函数在整个三轴区间广泛弥散，不存在静态固定的三轴取向。
\end{itemize}

这两种物理机制在低能唯象能谱上存在高度的简并性。例如，当 Davydov-Filippov 模型的刚性三轴角取为 $\gamma = 30^\circ$ 时，其预言的激发能级比 $E(2^+_2)/E(2^+_1) = 2.5$，这与 Wilets-Jean 理想 $\gamma$ 不稳定模型所预言的能级比完全一致。两类模型在低自旋四极跃迁分支比上也能通过调整微观参数产生高度退化的结果。因此，单凭能级比或少数几条能谱分支，根本无法区分核态究竟属于刚性三轴还是大振幅 $\gamma$ 软振动。

为了打破这种解释简并，第~2.1~节与第~2.4~节引入了模型无关的 Kumar-Cline 电四极矩不变量方法\cite{cline-1986-nuclear-shapes-coulomb-excitation}。该方法利用零角动量标量张量积在实验室参考系下的完备求和规则，将宏观电荷分布的内禀形变矩直接表达为多步库仑激发实验测得的一组约化矩阵元乘积：
\begin{equation}\label{eq:4-1-q2-invariant}
P_2 \equiv \langle Q^2 \rangle = \sqrt{5} \langle 0^+_i | [\mathcal{M}(E2) \times \mathcal{M}(E2)]_0 | 0^+_i \rangle = \sum_t \langle 0^+_i || \mathcal{M}(E2) || 2^+_t \rangle \langle 2^+_t || \mathcal{M}(E2) || 0^+_i \rangle \frac{(-1)^{-I_t}}{\sqrt{2I_i+1}},
\end{equation}
\begin{equation}\label{eq:4-1-q3-invariant}
P_3 \equiv \langle Q^3 \cos 3\delta \rangle = -\sqrt{\frac{35}{2}} \langle 0^+_i | [\mathcal{M}(E2) \times [\mathcal{M}(E2) \times \mathcal{M}(E2)]_2]_0 | 0^+_i \rangle.
\end{equation}
参数 $Q$ 对应宏观总四极形变大小，内禀角 $\delta$ 直接映射至三轴角。四阶不变量 $\langle Q^4 \rangle$ 则进一步给出了四极矩涨落与角向分散度 $\sigma(\delta) = \sqrt{\langle \delta^2 \rangle - \langle \delta \rangle^2}$ 的模型约束。

然而必须保持审慎的证据边界：低阶几何不变量 $\langle Q^2 \rangle$ 与 $\langle Q^3 \cos 3\delta \rangle$ 仅提供了宏观形状电荷分布的二阶与三阶矩重心，并不能唯一反演核多体波函数在 $(\beta,\gamma)$ 连续坐标空间上的全概率分布密度。更重要的是，实验测得的 $E2$ 矩阵元网络总是有限截断的，高能中间态的缺失与矩阵元相对符号的不确定性会引入系统误差。因此，分散度参数能够为检验多体模型波函数提供关键约束，但绝不能将其简化为某种机械数值门槛去绝对化地判定刚性或软核。

\subsection{实验探针与观测选择定则：自旋窗口、激发机制与电磁张量算符}\label{sec:4-1-4}

跨区比较结论的可信度，高度受制于所采用实验探针的布居物理机制及其自旋-宇称选择定则。不同实验手段在核素图上所能触及的激发能量与自旋窗口存在显著差异：
\begin{enumerate}
\item \textbf{低自旋、高精度绝对跃迁探针}：
   包括低能安全库仑激发（Coulomb Excitation）、$\beta$ 衰变谱学以及热中子/共振中子俘获与光核共振荧光（NRF）。这类探针的电磁相互作用或弱相互作用机制精确已知、微扰截面理论清晰，能够直接提取绝对约化跃迁几率 $B(E\lambda)$、$B(M\lambda)$ 以及对角静电四极矩 $Q_s$。但其自旋窗口通常受限于低角动量输入（多集中于 $I \le 3\text{--}4$），且对多步弱级联馈送极为敏感。
\item \textbf{中高自旋在束反应探针}：
   包括重离子熔合蒸发反应（Heavy-Ion Fusion-Evaporation）、深部非弹性碰撞（Deep Inelastic Collisions）、多核子转移反应以及自发/诱发裂变碎块谱学。这类探针通过大角动量输入能够有效布居高自旋转动带，是探索高自旋带交叉、磁转动剪刀带、手征双重带及摆动激发的有力工具。然而，必须注意不同反应机制的角动量输入截面与退激统计分布截然不同；且多步库仑激发在重区强形变核中同样可以达到 $I \sim 20\text{--}30$ 的高自旋态。绝对跃迁强度的测定不仅依赖反冲距离法（RDDS）和多普勒频移衰减法（DSAM），还广泛依赖于电子学超快计时（如 $\mathrm{LaBr}_3(\mathrm{Ce})$ 阵列）、延迟符合、转化电子谱学以及 NRF 共振荧光线宽分析等互补技术。
\end{enumerate}

此外，不同物理可观测量对核波函数的微观敏感度呈阶梯式分化：
\begin{itemize}
\item \textbf{能级能量 $E(I)$}（哈密顿量一阶期望值）：对多体相互作用微调较为迟钝，不同物理模型容易拟合出相近的谱形，判据敏感度最低；
\item \textbf{跃迁矩阵元 $B(E2)$ 与静四极矩 $Q_s$}（非对角与对角二阶电多极算符）：直接正比于内禀波函数的空间四极极化与电荷各向异性分布，是诊断转动模式与几何形状的核心基石；
\item \textbf{电单极跃迁强度 $\rho^2(E0)$ 与同位素位移 $\Delta\langle r^2 \rangle$}：严格依赖于两个不同态之间质子径向电荷密度分布的不均匀性重叠，对不同形变组态的微观混合具有最高灵敏度，是打破形状共存能谱简并的决定性利器；
\item \textbf{单核子转移反应与单粒子谱强}：实验直接测量的物理量是反应微分截面 $\mathrm{d}\sigma/\mathrm{d}\Omega$。提取单粒子光谱因子 $S_j$ 需要引入光学势畸变波波恩近似（DWBA）或绝热转移模型，其数值不仅依赖于反应动力学势的选择，还受限于实验激发能量分辨率以及截断窗口内的碎片求和；正如第~2.2~节所讨论的，单核子转移截面反映的是核多体波函数与单粒子试探算符的重叠，不能将其误解为裸算符的一体直接测量。
\end{itemize}

\subsection{可迁移与不可迁移条件的物理判据：六区比较总表}\label{sec:4-1-5}

基于上述微观机制与实验物理剖析，在跨质量区核结构比较中，必须严格区分普适可迁移的理论法则与不可迁移的局域条件：
\begin{itemize}
\item \textbf{可迁移的物理法则}：
  \begin{enumerate}
  \item 空间旋转与反演对称性的量子几何代数：Wigner-Eckart 定理的约化张量定义、Clebsch-Gordan 系数相位与 Racah 重耦合代数\cite{book-suhonen-from-nucleons-to-nucleus}；
  \item 电磁辐射场多极展开与跃迁几率的基本微观定义；
  \item 模型无关的 Kumar-Cline 四极不变量理论框架与零动量标量积完备求和规则\cite{cline-1986-nuclear-shapes-coulomb-excitation}；
  \item 态混合与微观组态干涉的代数分析方法。
  \end{enumerate}
\item \textbf{不可直接迁移的局域物理条件}：
  \begin{enumerate}
  \item 具体单粒子轨道的有效单粒子能量（ESPE）与局部主壳能隙；
  \item 价空间截断所引入的模型空间有效电荷（$e_p^{\mathrm{eff}}, e_n^{\mathrm{eff}}$）与有效弱轴矢流重整化因子 $g_A^{\mathrm{eff}}$；
  \item 唯象模型（如 IBM、PRM、玻尔哈密顿量 5DCH）通过局部能谱拟合所得的势能形状参数、微观质量张量或刚体转动惯量；
  \item 某一特定核区因特定高 $j$ 独特轨道单粒子耦合产生的复杂带结构（如手征双重带），在另一核区轨道占据与三轴极小尚未独立确立前，绝不能仅凭能谱外貌相似直接类推存在相同模式。
  \end{enumerate}
\end{itemize}

根据此项比较判据，表~\ref{tab:4-1-six-regions-framework}系统梳理了六个典型质量区的核心物理驱动、主导费米面轨道、关键形变特征以及当前实验认知的确定性边界。在全书后续比较论述中，我们将严格遵守表中关于``明确测定''、``未见''、``未搜索''与``灵敏度不足''四类认识论边界的界定，绝不以跨区类比掩盖物理认知的真实空缺。

\begin{table}[htbp]
\centering
\small
\caption{六个典型质量区结构演化比较的核心微观驱动、主导物理机制与实验认知边界}\label{tab:4-1-six-regions-framework}
\begin{tabularx}{\textwidth}{p{1.8cm}p{3.5cm}p{3.8cm}X}
\toprule
质量区 & 主导轨道与占据特征 & 核心物理机制与现象 & 实验证据与认知边界 \\
\midrule
$A \approx 60$ & $\pi,\nu(1f_{5/2}, 2p_{3/2}, 1g_{9/2})$ 弱跨壳 & 快速组态混合，次闭壳破缺，$N=Z$ 核区强质中配对 & $^{64}\mathrm{Ge}$ 测定寿命但静四极矩实测仍空缺；$^{68}\mathrm{Ni}$ 测得高压低能 $0^+$ 态与强 $E0$ 证实三形状共存 \\
\addlinespace
$A \approx 80$ & $\pi,\nu(1g_{9/2})$ 底部填充，强烈驱动形变 & 长椭与扁椭剧烈形状共存，大振幅 $\gamma$ 涨落与长扁反转 & $^{74,76}\mathrm{Kr}$ 多步库激测定静矩符号打破能谱简并；$^{76}\mathrm{Ge}$ 库激不变量证实刚性三轴特征；$^{74}\mathrm{As},^{78}\mathrm{Br}$ 奇奇核手征-反射待深入检验 \\
\addlinespace
$A \approx 100$ & $\pi(2p, 1f) \to 1g_{9/2}$ 跃迁，$\nu(1g_{7/2}, 1h_{11/2})$ 陡峭下移 & Type-II 壳演化诱发急剧形状量子相变，大振幅 $\gamma$ 软性，低能偶极与八极共振 & $^{98,100}\mathrm{Zr}$ 证实一阶相变；$^{98}\mathrm{Mo}$ 测得低能 $0_2^+$ 共存；$^{102}\mathrm{Ru}$ NRF 测得剪刀模与八极共振共存，绝非高自旋手征证据 \\
\addlinespace
$A \approx 130$ & $\pi(1g_{7/2}, 2d_{5/2}, 1h_{11/2})$ 与 $\nu(1h_{11/2})$ 顶部空穴 & 奇奇核与奇 $A$ 核手征双重带典型温床，横向与纵向摆动运动（wobbling）激辩 & 本论文核心锚点区；$^{135}\mathrm{Pr}, ^{133}\mathrm{Ce}, ^{131}\mathrm{Ba}$ 广泛报告三轴相关模式，但摆动解释面临 $\gamma$ 软势与多声子振动竞争 \\
\addlinespace
$A \approx 160$ & $\pi(1h_{11/2}), \nu(1i_{13/2})$ 强烈中部分裂，高能隙 & 典型大形变稳定长椭球转动，高 $K$ 异能态，微观配对减弱与准粒子排列 & 经典刚性转动基准区；Lu 同位素报告摆动激发，稀土高 $K$ 异能态展现极高受阻跃迁度 \\
\addlinespace
$A \approx 190$ & 费米面逼近 $Z=82, N=126$ 闭壳，$\nu(1i_{13/2})$ 空穴主导 & 稳定扁椭形变，$\gamma$ 软性与手征转动向球形单粒子壳模型的跨越过渡 & 测得扁椭转动带；Ir, Os, Au 同位素存在三轴特征但面临强八极与复合成核路径限制 \\
\bottomrule
\end{tabularx}
\end{table}
